\documentclass[11pt]{article}
\def\activityid{F08-F-v3}\usepackage[letterpaper,margin=.65in,top=.60in,bottom=.65in]{geometry}
\usepackage[T1]{fontenc}
\usepackage[default]{sourcesanspro}
\usepackage{microtype,tikz,array,tabularx,fancyhdr,amsmath,amssymb}
\usepackage[hidelinks]{hyperref}
\usetikzlibrary{calc,arrows.meta}
\setlength{\parindent}{0pt}\setlength{\parskip}{7pt}
\setlength{\headheight}{0pt}\setlength{\footskip}{26pt}\setlength{\emergencystretch}{2em}
\pagestyle{fancy}\fancyhf{}\renewcommand{\headrulewidth}{0pt}\renewcommand{\footrulewidth}{.4pt}
\fancyfoot[L]{\scriptsize Bellingham Math Circle / Week 8 / \activityid}
\fancyfoot[R]{\scriptsize \thepage}
\definecolor{ink}{gray}{.12}\definecolor{soft}{gray}{.95}\color{ink}
\newcommand{\PageTitle}[2]{{\fontsize{10}{12}\selectfont\bfseries WEEK 8\quad STRATEGY LAB\hfill #1\par}\vspace{3pt}{\fontsize{26}{29}\selectfont\bfseries #2\par}\vspace{2pt}}
\newcommand{\Subhead}[1]{\par\vspace{3pt}{\large\bfseries #1}\par}
\newcommand{\RuleBox}[1]{\par\begingroup\setlength{\fboxsep}{8pt}\colorbox{soft}{\parbox{\dimexpr\linewidth-16pt\relax}{#1}}\endgroup\par}
\newcommand{\WriteBox}[2]{\par\textbf{#1}\par\vspace{2pt}\begin{tikzpicture}\draw[black!40,line width=.5pt] (0,0) rectangle (\dimexpr\linewidth-.5pt\relax,#2);\end{tikzpicture}\par}
\newcommand{\RookBoard}[3]{\begin{tikzpicture}[x=#3,y=#3]\draw[step=1,black!45] (0,0) grid (#1,#2);\draw[thick] (0,0) rectangle (#1,#2);\node[font=\Large] at ({#1-.5},.5){$\star$};\draw[-{Latex},thick] (0,-.35)--(#1,-.35);\draw[-{Latex},thick] ({#1+.35},#2)--({#1+.35},0);\end{tikzpicture}}
\newcommand{\Table}[3]{\begin{tikzpicture}[x=#3,y=#3]\draw[step=1,black!25] (0,0) grid (#1,#2);\draw[thick] (0,0) rectangle (#1,#2);\fill (0,0)circle(3pt);\draw[-{Latex},very thick] (0,0)--(.7,.7);\node[below] at ({#1/2},-.1){width #1};\node[rotate=90,above] at (-.2,{#2/2}){height #2};\end{tikzpicture}}
\newcommand{\GraphNodes}{\coordinate(A)at(0,0);\coordinate(B)at(3,0);\coordinate(C)at(1.5,2.6);\coordinate(D)at(1.5,.95);}
\newcommand{\Kfour}[1]{\begin{tikzpicture}[scale=#1]\GraphNodes\draw[very thick](A)--(B)--(C)--(A)--(D)--(B) (D)--(C);\foreach \n in {A,B,C,D}{\node[circle,draw,fill=white,minimum size=22pt,font=\bfseries]at(\n){\n};}\end{tikzpicture}}
\newcommand{\House}[1]{\begin{tikzpicture}[scale=#1]\coordinate(A)at(0,0);\coordinate(B)at(3,0);\coordinate(C)at(3,2);\coordinate(D)at(0,2);\coordinate(E)at(1.5,3.3);\draw[very thick](A)--(B)--(C)--(D)--(A) (D)--(E)--(C);\foreach \n in {A,B,C,D,E}{\node[circle,draw,fill=white,minimum size=22pt,font=\bfseries]at(\n){\n};}\end{tikzpicture}}




\newcommand{\StudentPage}[1]{{\fontsize{10}{12}\selectfont\bfseries Week 8 / Strategy lab / #1\par}\vspace{10pt}}
\newcommand{\Problem}[2]{\par\noindent\textbf{Problem #1:} #2\par}
\newcommand{\Space}[1]{\par\begin{tikzpicture}\draw[black!35] (0,0) rectangle (\dimexpr\linewidth-.5pt\relax,#1);\end{tikzpicture}\par}
\newcommand{\Piles}[1]{\begin{center}\begin{tikzpicture}[x=1in,y=1in]\foreach \x in {0,...,#1}{\draw[black!50,rounded corners] ({\x*2.05},0) rectangle ({\x*2.05+1.8},1.3);}\end{tikzpicture}\end{center}}

\begin{document}

\PageTitle{FACILITATOR / 1 OF 6}{From play to a useful record}
\RuleBox{F08 v3, September 2026: prepared and unpiloted. These changes apply the organizer's Week 1 feedback; there is no reported Week 8 classroom outcome. The goal is a strategy discovered through repeated concrete games. Pages are a menu for more than one possible meeting.}
\Subhead{Prepare a small starting set}
Print page 1 for each child at their starting level: three K--1, three middle, one upper (7 sheets). Keep one master of each continuation and the 3-page extra; copy only what is needed. Middle page 3 reuses its earlier board. Bring about 60 counters, five plates or marked pile areas, pencils, and scraps labeled 1,2,4,8. Printed pile mats are usable counter areas; reset and reuse them.

K--1: adult reads; compare distances/count to 3. Middle: count to 7 and track moves, with adult scribing if needed. Upper: small subtraction and odd/even; bundles introduce binary notation. Extra: organize lists and compare differences; use a calculator only for the final supplied formula. Choose by prerequisites, not grade.
\Subhead{Common launch before separate starts}
0--10: let everyone move counters and build piles. 10--15: gather at one visible rook board. Show a one-square right move, a longer down move, and the same token reaching the star. Invite a child to make a legal move; have another mark its start and finish. Say: ``One token, right or down, at least one square; reaching the star wins.'' Show an illegal diagonal only if needed. Do not reveal equal distances yet. Place the token at (2,1) and build matching piles of 2 and 1. Slide right one square while removing one counter from the 2-pile. State that pile games remove any positive number from just one pile; taking the last counter wins. Then distribute starting tasks. This one correspondence shows both actions without revealing a strategy.

15--35: play the contrasting starts. 35--40: movement/reset. 40--55: repeat, compare, or introduce the next record only when it answers the children's question. 55--60: share one move with its reason and tidy. The timetable is an adaptable proposal, not a source-prescribed duration.

Parent stays mainly with K--1. Organizer alternates middle/upper about every five minutes, leaving a specific next game. Triplets rotate two opponents and a checker. The upper child needs actual games and discussion with an adult.
\Subhead{Pacing, hints, and stopping}
Ask for a legal move before a strategy. Then ask what is left to the star, or compare an equal and an unequal start. K can stop after showing the two center moves. Middle can stop with an argued board pattern or the board/pile correspondence. Upper can stop after the six replies to (1,2,3), or after balancing concrete bundle records. No general proof is a printed completion requirement; pages 2--3 retain the proofs for conversation. The extra can stop after its small map.
\newpage

\PageTitle{FACILITATOR / 2 OF 6}{Checked two-pile reasoning}
\Subhead{K--1 and middle board answers}
Distances $(a,b)$ mean squares remaining right and down. A rook move reduces exactly one positive coordinate. The star is $(0,0)$: the previous player has already won. On the $3\times3$ board, the center $(1,1)$ is losing. Both possible moves allow an immediate reply to the star. Top-left $(2,2)$ is also losing; copy the amount removed in the other direction.

On the $4\times4$ board, the traps form the top-left to bottom-right diagonal: $(3,3),(2,2),(1,1),(0,0)$. Every unequal pair is winning: reduce the larger coordinate to the smaller. The middle start $(3,2)$ has the unique winning move to $(2,2)$.
\Subhead{A complete proof, with counters}
If two piles have the same size, every move makes them unequal. If they differ, reducing the larger to the smaller is legal and restores equality. Thus a player who hands over equal piles can always restore equality after an opponent's move. The total number of counters strictly decreases, so play must finish. The player restoring equality takes the final counters; the opponent receives $(0,0)$ and cannot move. These two properties establish all traps, on rectangles as well as squares.

Any unequal two-pile start has exactly one winning first move (specifying the pile and amount). Matching piles have none. Children can use $(5,2)$ as their designed puzzle: remove 3 from the 5-pile.
\Subhead{Why three piles change the problem}
$(1,1,1)$ is winning: remove one pile and leave $(0,1,1)$. Even total is insufficient: $(1,1,2)$ has total 4 but is winning by removing the whole 2-pile. $(1,2,3)$ is losing. Here are all first moves and one reply; pile order is kept fixed.
\begin{center}\renewcommand{\arraystretch}{1.35}\begin{tabular}{c|c@{\qquad}c|c}After move&Reply leaves&After move&Reply leaves\\\hline(0,2,3)&(0,2,2)&(1,2,0)&(1,1,0)\\(1,0,3)&(1,0,1)&(1,2,1)&(1,0,1)\\(1,1,3)&(1,1,0)&(1,2,2)&(0,2,2)\end{tabular}\end{center}
Each reply leaves two equal piles and one empty pile. This gives a physical proof before binary notation. The middle group can stop here or join the upper bundle experiment.

\newpage

\PageTitle{FACILITATOR / 3 OF 6}{The binary invariant}
\Subhead{Upper answers}
Binary records: $1=001$, $2=010$, $3=011$, $4=100$, $5=101$, $6=110$, $7=111$. $(1,2,3)$ and $(2,4,6)$ have even column sums. $(1,2,4)$ does not; reduce 4 to 3. For the table: $(2,4,6)$ is balanced; $(3,4,5)$ can become $(1,4,5)$; $(1,5,7)$ can become $(1,5,4)$; $(3,5,6)$ is balanced. Other correct balancing moves should be accepted.

For $(7,10,12)$, the Nim sum is 1; reduce 7 to 6. For chosen first piles $a,b$, the unique third pile is $a\mathbin{\oplus}b$: each binary digit is forced by whether the first two digits agree.
\Subhead{Why the test is sufficient, for any number of piles}
A move changes exactly one integer. At least one of its binary digits changes, so at least one previously even column becomes odd. No legal move can connect two balanced positions.

If the Nim sum $s$ is nonzero, choose its highest 1-bit, say the $2^k$ column. Some pile $x$ has a 1 there. Replace it by $y=x\mathbin{\oplus}s$. Higher bits stay fixed, this bit becomes 0, and lower bits may change. The sum of all smaller bundles is $2^k-1$, so $y<x$: the move is legal. Every odd column is toggled once, every even column is unchanged, and the result is balanced. The decreasing total ensures termination. These two facts prove exactly the same win/lose structure as in the two-pile proof.

Do not physically permit removing an arbitrary bundle while leaving a larger number of counters: the game is still ordinary removal from one pile. The bundle display is a way to calculate a legal smaller result.
\Subhead{Extra: Wythoff, checked results and proof boundary}
The old equal-pile trap fails because both piles can be emptied in one diagonal move. For $0\le a\le b\le7$, the traps are $(0,0),(1,2),(3,5),(4,7)$. Greedy pairs through $n=5$ are $(0,0),(1,2),(3,5),(4,7),(6,10),(8,13)$.

Distinct pairs have disjoint coordinates and distinct absolute differences, so neither a single-pile move nor an equal-take move connects them, including reversed pile order. This proves only the no-move-between-traps condition. To prove completeness one must also show every outside position reaches a listed pair. Sample moves: $(2,4)\to(2,1)$, $(5,8)\to(4,7)$, $(8,12)\to(6,10)$. The golden-ratio formula is a cited theorem, not a conclusion from five data points.

\newpage

\PageTitle{FACILITATOR / 4 OF 6}{Where the mathematics leads}
\Subhead{Undergraduate direction and genuine research}
Two piles give a strategy-preserving change of representation from a grid to an impartial game. Binary column parity is addition in a vector space over the two-element field. The proof has a useful algorithmic form: recognize the losing set, prove no move stays inside it, and construct a move into it from everywhere else. Formal algebra vocabulary can wait.

Charles Bouton's historical paper, \emph{Nim, a Game with a Complete Mathematical Theory}, \emph{Annals of Mathematics} 3 (1901--02), pp. 35--39, is the original research source for binary analysis. Our normal-play proof is written out on page 3, so the circle does not depend on reading the historical notation. \href{https://paradise.caltech.edu/ist4/lectures/Bouton1901.pdf}{Caltech-hosted paper}.

Fraenkel and Peled, \emph{Harnessing the unwieldy MEX function}, \emph{Games of No Chance 4}, MSRI Publications 63 (2015), pp. 77--94: \S1, pp. 77--78 states the Wythoff rule, greedy pairs, and golden-ratio formula; \S5 develops efficient computation for generalized complementary sequences. \href{https://library.slmath.org/books/Book63/files/131104-Fraenkel-2.pdf}{Author research chapter}.

The extra packet meets an actual research theme: a simple rule change can replace bitwise arithmetic with complementary sequences and a question of efficient computation. It does not reproduce the paper's generalized algorithm or claim an open problem has been solved by the classroom experiment.
\Subhead{Local sources and teaching choices}
JRMF, \emph{Rook's Move Activity Guide}, PDF pp. 3--6: legal moves, exploration before advice, losing positions, and the Wythoff connection. We replace the old upper king variant with three-pile Nim; the king version remains an untaught reserve unless the use log says otherwise.

\emph{Math Circle by the Bay}, preface printed pp. viii--x (PDF pp. 9--11), supports common themes with adjustable depth, manipulatives, explanations, and extra challenges. Rozhkovskaya, Lesson 7 ``At the lesson,'' EPUB \texttt{part0017.xhtml}, highlights checking that children understand legal moves. The seven-child schedule and bundle prompts are our adaptations.

Lowell Handout 2, problem 2.3, and Handout 3, problems 3.1--3.3, include subtraction-game history; Week 7 also studies constrained subtraction. Here an entire pile may be removed, and the important new step is multiple piles plus binary balance. Reuse is deliberate, not evidence that every child has seen these exact starts.
\Subhead{Record and verify}
IDs F08-K/M/U/X-v3. Record the actual starts, rule, level, strategies explained, and whether Wythoff was attempted. Keep unused stages for later years. \texttt{verify.py} checks all three-pile states through size 12 and Wythoff states through 30 by independent backward recursion. Those finite checks support the examples; the proofs establish the unlimited claims.


\newpage
\PageTitle{FACILITATOR / 5 OF 6}{Use each record for a reason}
\Subhead{Core sequence and specific hint ladders}
\textbf{K pages 1--2:} compare an immediate win with the center, then equal top-left distances with an unequal start. Demonstrate drawing a move from square center to square center. When stuck, enumerate the center's two legal moves with the token. Page 3 is a separate translation: build two equal piles, copy a removal, then change to unequal piles. Stop before this transition if rook play remains engaging.

\textbf{Middle pages 1--2:} play eight purposeful starts before the T/W map. Demonstrate the (right-distance, down-distance) record at (2,1) with the actual token. The first table records game outcomes, not proved winning types. Before mapping, classify (1,0) and (1,1) together: a W has one move to T, a T has every move to W. Ask children to challenge one proposed letter; a lost game alone proves neither T nor W. Page 3 physically matches two slides to removals before asking for equal-pile replies. Page 4 adds three piles only after equal-pile play is understood. The six preprinted outcomes prevent copying from consuming the investigation.

\textbf{Upper pages 1--2:} compare adjacent starts and test the even-total guess. Two equal piles plus an empty pile is useful, but (2,4,6) has no immediate move to that form: failing to find one is a reason to seek another representation. The six replies to (1,2,3) can be a complete stopping point before binary. Page 3 groups real counters before recording 0/1 columns. Model 5 as 4+1, then let the child build. If notation obscures the game, leave bundle cards visible beside ordinary piles. Page 4 tests the same record on four cases and actual opponent moves. Page 5 introduces the highest-odd-column recipe only after the repair problem is familiar.

\textbf{Proof reserve:} ask why a move from balance must alter a column; then why the highest changed bundle outweighs all lower bundles together; finally use decreasing total for termination. Preserve both directions of the losing-set argument. Naming XOR is optional and follows the child's record, not the launch.
\Subhead{Extra and observation record}
Play equal and unequal starts before the classification grid. Work only far enough to reveal a second nontrivial trap; finish the grid later if useful. Demonstrate sorting after a move: (2,4) can become (2,1), recorded at (1,2). The greedy sequence is a compressed record of already encountered traps. Distinct coordinates/differences prove no internal moves; completeness is a separate requirement. Supply the golden-ratio result as a theorem only if reached.

After teaching, record exact starts/pages, who needed a rule repeated, useful or burdensome records, child conjectures, and requests to continue. Separate observations from hypotheses and adult-supplied theorems. Mark unused pages unpiloted; do not infer uptake from printing.
\newpage
\PageTitle{FACILITATOR / 6 OF 6}{Additional v3 solution checks}
\Subhead{Board and pile instances}
K1 immediate-above-star is (0,1), winning immediately; the center is (1,1), losing. From top-left (2,2), right 1 is answered down 1; right 2 by down 2, and vice versa. The middle-left start (2,1) can move right 1 to (1,1). Equal pile starts (2,2),(3,3) use matching replies; (3,1) removes 2 from the larger pile.

Middle trial starts (1,0),(2,1),(3,1),(3,2),(2,3) are W; (1,1),(2,2),(3,3) are T. Unique moves to equal distances respectively leave (0,0),(1,1),(1,1),(2,2),(2,2). At (3,2), right 1 corresponds to piles (2,2); down 2 to (3,0). In pile Problem 6, (3,2) reduces to (2,2), (1,4) to (1,1); equal starts require an opponent's move first. The six replies in middle Problem 8 and upper Problem 4 are listed on page 2.

\Subhead{Upper bundle work}
Problem 3: (1,1,1) removes any single 1; (1,1,2) removes the 2. Neither (1,2,3) nor (2,4,6) reaches two equal nonempty piles and an empty pile in one move, since no pair is initially equal. Their totals 6 and 12 are even, but the total-4 winning example (1,1,2) refutes even-total sufficiency.

Problem 5 records in 4/2/1 order are 1=001, 2=010, 3=011, 4=100, 5=101, 6=110, 7=111. Problem 6 column counts for (1,2,3) are (0,2,2); for (1,2,4), (1,1,1). Reduce 4 to 3. Problem 8 checked moves are on page 3; multiple balancing moves can occur.

Problem 10: (4,6,7) has odd 4 and 1 columns. Legal repairs are 4 to 1, 6 to 3, or 7 to 2. In (7,10,12), only the 1 column is odd; reduce 7 to 6. Problem 11 third piles are 7 for (2,5) and 5 for (3,6); each bit is forced. Accept a child-designed zero-XOR start with a demonstrated legal reply.
\Subhead{Extra instances}
Problem 1: (2,2) and (2,3) are winning, moving respectively to (0,0) and (1,2); (1,2),(3,5) are traps. Problem 2's four outcomes are (0,2),(1,1),(1,0),(0,1); empty the remaining pile or take 1 from both. Problem 3 grid has exactly (0,0),(1,2),(3,5),(4,7) as T. Problem 4: (2,4) to (2,1); (4,6) to (3,5); (5,7) to (3,5), all legal. Greedy/formula rows and final large moves are on page 3. For each tested pair, single-pile removal would need a shared coordinate; equal removal would need the same difference. Neither occurs between distinct listed traps.

\end{document}
